% --- Archon physics formalization source begin ---
% archon:physics
% archon:covers PhyXMiniProblems/problem_phyx_mini_0617.lean
% archon:source-report reports/phyx_mini/problem_phyx_mini_0617.source.json

\chapter{Physics problem phyx\_mini\_0617}
\label{ch:PhyXMiniProblems_problem_phyx_mini_0617}

\paragraph{Problem source.}
\#\# Physical scenario

Consider a potential well defined as  \$U(x) = \textbackslash{}infty\$ for \$x < 0\$, \$U(x) = 0\$ for \$0 < x < L\$, and \$U(x) = U\_0 > 0\$ for \$x > L\$ (figure). Consider a particle with mass \$m\$ and kinetic energy \$E < U\_0\$ that is trapped in the well.The wave function must remain finite as \$x \textbackslash{}to \textbackslash{}infty\$.And it satisfy both the Schrödinger equation and this boundary condition at infinity.

\#\# Question

What must the form of the function \$\textbackslash{}psi(x)\$ for \$x > L\$ be?

\#\# Answer choices

- A: A: \$Ce\textasciicircum{}\{\textbackslash{}kappa x\} + De\textasciicircum{}\{-\textbackslash{}kappa x\}\$

- B: B: \$Ae\textasciicircum{}\{\textbackslash{}kappa x\} + Be\textasciicircum{}\{-\textbackslash{}kappa x\}\$

- C: C: \$Fe\textasciicircum{}\{\textbackslash{}kappa x\} + Ge\textasciicircum{}\{-\textbackslash{}kappa x\}\$

- D: D: \$Xe\textasciicircum{}\{\textbackslash{}kappa x\} + Ye\textasciicircum{}\{-\textbackslash{}kappa x\}\$

\#\# Figure caption (auxiliary; use the image as primary evidence)

The image is a graph depicting a potential energy function \textbackslash{}( U(x) \textbackslash{}) in relation to the position \textbackslash{}( x \textbackslash{}). Here's a detailed description:

- **Axes:**
  - The vertical axis is labeled \textbackslash{}( U(x) \textbackslash{}) and represents the potential energy. It has a value marked at infinity (\textbackslash{}(\textbackslash{}infty\textbackslash{})) and another point labeled \textbackslash{}( U\_0 \textbackslash{}).
  - The horizontal axis is labeled \textbackslash{}( x \textbackslash{}) and represents position.

- **Graph Characteristics:**
  - The potential starts at infinity at \textbackslash{}( x = 0 \textbackslash{}) and instantly drops to zero.
  - From \textbackslash{}( x = 0 \textbackslash{}) to \textbackslash{}( x = L \textbackslash{}), the potential \textbackslash{}( U(x) \textbackslash{}) is zero, forming a horizontal line along the x-axis.
  - At \textbackslash{}( x = L \textbackslash{}), there is a sudden increase to \textbackslash{}( U\_0 \textbackslash{}), marked by a vertical line.
  - From \textbackslash{}( x = L \textbackslash{}) onward, the potential remains constant at \textbackslash{}( U\_0 \textbackslash{}), forming a horizontal line parallel to the x-axis.

- **Shading and Lines:**
  - The potential barrier at the start (\textbackslash{}( x = 0 \textbackslash{})) and end (\textbackslash{}( x = L \textbackslash{})) is highlighted with shading.
  - A horizontal dashed line across the graph indicates the level \textbackslash{}( U\_0 \textbackslash{}). 

This setup resembles a finite potential well, which is a common concept in quantum mechanics.

\#\# Dataset metadata

Quantum Phenomena; Physical Model Grounding Reasoning, Numerical Reasoning

\paragraph{Recorded answer/context.}
A: A: \$Ce\textasciicircum{}\{\textbackslash{}kappa x\} + De\textasciicircum{}\{-\textbackslash{}kappa x\}\$

\paragraph{Figure/image path.}
/root/proposal\_for\_physic/science-mango/phyx\_mini\_run/phyx\_data/test\_image/617.png

\paragraph{Formalization target.}
create a compiling Lean file with sorry bodies at `PhyXMiniProblems/problem\_phyx\_mini\_0617.lean`. The Lean declarations must preserve the physical quantities, dimensions or dimensional roles, figure labels, governing-law hypotheses, and final relation expressed by this problem.
Use Mathlib/Physlib names found through LeanExplore where available. If a physics API is missing, introduce faithful local abstractions rather than scalar placeholder aliases.

\begin{theorem}[Physics formalization target]
\label{thm:physics:phyx_mini_0617:target}
\lean{PhyXMiniProblems.ProblemPhyXMini0617.exterior_waveFunction_has_choice_A_form}
\uses{def:physics:phyx-mini-0617:phyxminiproblems-problemphyxmini0617-wavenumberquantity, def:physics:phyx-mini-0617:phyxminiproblems-problemphyxmini0617-lengthinmeters, def:physics:phyx-mini-0617:phyxminiproblems-problemphyxmini0617-massinkilograms, def:physics:phyx-mini-0617:phyxminiproblems-problemphyxmini0617-energyinjoules, def:physics:phyx-mini-0617:phyxminiproblems-problemphyxmini0617-actioninjouleseconds, def:physics:phyx-mini-0617:phyxminiproblems-problemphyxmini0617-wavenumberininversemeters, def:physics:phyx-mini-0617:phyxminiproblems-problemphyxmini0617-finitepotentialwellsetup, def:physics:phyx-mini-0617:phyxminiproblems-problemphyxmini0617-matchespotentialwellscenario, def:physics:phyx-mini-0617:phyxminiproblems-problemphyxmini0617-matchessuppliedpotentialwellfigure, def:physics:phyx-mini-0617:phyxminiproblems-problemphyxmini0617-hasphysicalpotentialwellparameters, def:physics:phyx-mini-0617:phyxminiproblems-problemphyxmini0617-usesstandardreducedplanckconstant, def:physics:phyx-mini-0617:phyxminiproblems-problemphyxmini0617-satisfiesexteriorstationaryschrodingerequation, def:physics:phyx-mini-0617:phyxminiproblems-problemphyxmini0617-satisfiesfiniteboundaryconditionatinfinity}
The assigned autoformalize agent should translate this physics problem into Lean declarations in the covered file, with theorem and lemma proof bodies written as `by sorry`.
\end{theorem}
\begin{proof}
This is an autoformalization task, not a proof task. Produce statements that can later be proved by the physics prover without weakening the physical model.
\end{proof}

% --- Archon declaration topology begin ---
\section{Declaration topology}
\begin{definition}[Length Quantity]
  \label{def:physics:phyx-mini-0617:phyxminiproblems-problemphyxmini0617-lengthquantity}
  \lean{PhyXMiniProblems.ProblemPhyXMini0617.LengthQuantity}
  A nonnegative physical length, independent of the choice of units
\end{definition}
\begin{proof}
  This is the declared model definition of Length Quantity.
\end{proof}
\begin{definition}[Mass Quantity]
  \label{def:physics:phyx-mini-0617:phyxminiproblems-problemphyxmini0617-massquantity}
  \lean{PhyXMiniProblems.ProblemPhyXMini0617.MassQuantity}
  A nonnegative physical mass, independent of the choice of units
\end{definition}
\begin{proof}
  This is the declared model definition of Mass Quantity.
\end{proof}
\begin{definition}[action Dimension]
  \label{def:physics:phyx-mini-0617:phyxminiproblems-problemphyxmini0617-actiondimension}
  \lean{PhyXMiniProblems.ProblemPhyXMini0617.actionDimension}
  The physical dimension of action, mass * length² / time
\end{definition}
\begin{proof}
  This is the declared model definition of action Dimension.
\end{proof}
\begin{definition}[Action Quantity]
  \label{def:physics:phyx-mini-0617:phyxminiproblems-problemphyxmini0617-actionquantity}
  \lean{PhyXMiniProblems.ProblemPhyXMini0617.ActionQuantity}
  \uses{def:physics:phyx-mini-0617:phyxminiproblems-problemphyxmini0617-actiondimension}
  A nonnegative physical action, used for the reduced Planck constant
\end{definition}
\begin{proof}
  This is the declared model definition of Action Quantity.
\end{proof}
\begin{definition}[Wave Number Quantity]
  \label{def:physics:phyx-mini-0617:phyxminiproblems-problemphyxmini0617-wavenumberquantity}
  \lean{PhyXMiniProblems.ProblemPhyXMini0617.WaveNumberQuantity}
  A nonnegative inverse length, the dimension carried by κ
\end{definition}
\begin{proof}
  This is the declared model definition of Wave Number Quantity.
\end{proof}
\begin{definition}[length In Meters]
  \label{def:physics:phyx-mini-0617:phyxminiproblems-problemphyxmini0617-lengthinmeters}
  \lean{PhyXMiniProblems.ProblemPhyXMini0617.lengthInMeters}
  \uses{def:physics:phyx-mini-0617:phyxminiproblems-problemphyxmini0617-lengthquantity}
  Metre readout of a physical length
\end{definition}
\begin{proof}
  This is the declared model definition of length In Meters.
\end{proof}
\begin{definition}[mass In Kilograms]
  \label{def:physics:phyx-mini-0617:phyxminiproblems-problemphyxmini0617-massinkilograms}
  \lean{PhyXMiniProblems.ProblemPhyXMini0617.massInKilograms}
  \uses{def:physics:phyx-mini-0617:phyxminiproblems-problemphyxmini0617-massquantity}
  Kilogram readout of a physical mass
\end{definition}
\begin{proof}
  This is the declared model definition of mass In Kilograms.
\end{proof}
\begin{definition}[energy In Joules]
  \label{def:physics:phyx-mini-0617:phyxminiproblems-problemphyxmini0617-energyinjoules}
  \lean{PhyXMiniProblems.ProblemPhyXMini0617.energyInJoules}
  Joule readout of a physical energy
\end{definition}
\begin{proof}
  This is the declared model definition of energy In Joules.
\end{proof}
\begin{definition}[action In Joule Seconds]
  \label{def:physics:phyx-mini-0617:phyxminiproblems-problemphyxmini0617-actioninjouleseconds}
  \lean{PhyXMiniProblems.ProblemPhyXMini0617.actionInJouleSeconds}
  \uses{def:physics:phyx-mini-0617:phyxminiproblems-problemphyxmini0617-actionquantity}
  Joule-second readout of a physical action
\end{definition}
\begin{proof}
  This is the declared model definition of action In Joule Seconds.
\end{proof}
\begin{definition}[wave Number In Inverse Meters]
  \label{def:physics:phyx-mini-0617:phyxminiproblems-problemphyxmini0617-wavenumberininversemeters}
  \lean{PhyXMiniProblems.ProblemPhyXMini0617.waveNumberInInverseMeters}
  \uses{def:physics:phyx-mini-0617:phyxminiproblems-problemphyxmini0617-wavenumberquantity}
  Inverse-metre readout of a physical wave number
\end{definition}
\begin{proof}
  This is the declared model definition of wave Number In Inverse Meters.
\end{proof}
\begin{definition}[zero Potential Energy]
  \label{def:physics:phyx-mini-0617:phyxminiproblems-problemphyxmini0617-zeropotentialenergy}
  \lean{PhyXMiniProblems.ProblemPhyXMini0617.zeroPotentialEnergy}
  The unit-independent physical zero of potential energy
\end{definition}
\begin{proof}
  This is the declared model definition of zero Potential Energy.
\end{proof}
\begin{definition}[Potential Value]
  \label{def:physics:phyx-mini-0617:phyxminiproblems-problemphyxmini0617-potentialvalue}
  \lean{PhyXMiniProblems.ProblemPhyXMini0617.PotentialValue}
  A point of the potential is either a finite energy or an infinite wall
\end{definition}
\begin{proof}
  This is the declared model definition of Potential Value.
\end{proof}
\begin{definition}[Figure Label]
  \label{def:physics:phyx-mini-0617:phyxminiproblems-problemphyxmini0617-figurelabel}
  \lean{PhyXMiniProblems.ProblemPhyXMini0617.FigureLabel}
  Literal labels visible in the supplied potential-energy graph
\end{definition}
\begin{proof}
  This is the declared model definition of Figure Label.
\end{proof}
\begin{definition}[Supplied Potential Well Figure]
  \label{def:physics:phyx-mini-0617:phyxminiproblems-problemphyxmini0617-suppliedpotentialwellfigure}
  \lean{PhyXMiniProblems.ProblemPhyXMini0617.SuppliedPotentialWellFigure}
  \uses{def:physics:phyx-mini-0617:phyxminiproblems-problemphyxmini0617-figurelabel}
  Qualitative and textual readouts from the supplied raster image
\end{definition}
\begin{proof}
  This is the declared model definition of Supplied Potential Well Figure.
\end{proof}
\begin{definition}[Finite Potential Well Setup]
  \label{def:physics:phyx-mini-0617:phyxminiproblems-problemphyxmini0617-finitepotentialwellsetup}
  \lean{PhyXMiniProblems.ProblemPhyXMini0617.FinitePotentialWellSetup}
  \uses{def:physics:phyx-mini-0617:phyxminiproblems-problemphyxmini0617-lengthquantity, def:physics:phyx-mini-0617:phyxminiproblems-problemphyxmini0617-massquantity, def:physics:phyx-mini-0617:phyxminiproblems-problemphyxmini0617-actionquantity, def:physics:phyx-mini-0617:phyxminiproblems-problemphyxmini0617-potentialvalue, def:physics:phyx-mini-0617:phyxminiproblems-problemphyxmini0617-suppliedpotentialwellfigure}
  The independent physical data. The source's kinetic energy E is measured in the zero-potential interior and hence is also the stationary-state energy. The wave function has complex amplitude, while its argument is a metre readout
\end{definition}
\begin{proof}
  This is the declared model definition of Finite Potential Well Setup.
\end{proof}
\begin{definition}[Matches Potential Well Scenario]
  \label{def:physics:phyx-mini-0617:phyxminiproblems-problemphyxmini0617-matchespotentialwellscenario}
  \lean{PhyXMiniProblems.ProblemPhyXMini0617.MatchesPotentialWellScenario}
  \uses{def:physics:phyx-mini-0617:phyxminiproblems-problemphyxmini0617-lengthinmeters, def:physics:phyx-mini-0617:phyxminiproblems-problemphyxmini0617-zeropotentialenergy, def:physics:phyx-mini-0617:phyxminiproblems-problemphyxmini0617-finitepotentialwellsetup}
  The piecewise potential stated in the problem. This premise contains no exponential ansatz and no information about a coefficient or decay rate
\end{definition}
\begin{proof}
  This is the declared model definition of Matches Potential Well Scenario.
\end{proof}
\begin{definition}[Matches Supplied Potential Well Figure]
  \label{def:physics:phyx-mini-0617:phyxminiproblems-problemphyxmini0617-matchessuppliedpotentialwellfigure}
  \lean{PhyXMiniProblems.ProblemPhyXMini0617.MatchesSuppliedPotentialWellFigure}
  \uses{def:physics:phyx-mini-0617:phyxminiproblems-problemphyxmini0617-finitepotentialwellsetup}
  Text and qualitative facts read directly from image 617.png
\end{definition}
\begin{proof}
  This is the declared model definition of Matches Supplied Potential Well Figure.
\end{proof}
\begin{definition}[Has Physical Potential Well Parameters]
  \label{def:physics:phyx-mini-0617:phyxminiproblems-problemphyxmini0617-hasphysicalpotentialwellparameters}
  \lean{PhyXMiniProblems.ProblemPhyXMini0617.HasPhysicalPotentialWellParameters}
  \uses{def:physics:phyx-mini-0617:phyxminiproblems-problemphyxmini0617-lengthinmeters, def:physics:phyx-mini-0617:phyxminiproblems-problemphyxmini0617-massinkilograms, def:physics:phyx-mini-0617:phyxminiproblems-problemphyxmini0617-energyinjoules, def:physics:phyx-mini-0617:phyxminiproblems-problemphyxmini0617-actioninjouleseconds, def:physics:phyx-mini-0617:phyxminiproblems-problemphyxmini0617-finitepotentialwellsetup}
  The source conditions m > 0, U₀ > 0, and 0 ≤ E < U₀, together with the positivity of L and ℏ. These are parameter facts rather than any part of the requested exterior wave-function formula
\end{definition}
\begin{proof}
  This is the declared model definition of Has Physical Potential Well Parameters.
\end{proof}
\begin{definition}[Uses Standard Reduced Planck Constant]
  \label{def:physics:phyx-mini-0617:phyxminiproblems-problemphyxmini0617-usesstandardreducedplanckconstant}
  \lean{PhyXMiniProblems.ProblemPhyXMini0617.UsesStandardReducedPlanckConstant}
  \uses{def:physics:phyx-mini-0617:phyxminiproblems-problemphyxmini0617-actioninjouleseconds, def:physics:phyx-mini-0617:phyxminiproblems-problemphyxmini0617-finitepotentialwellsetup}
  Calibration of the setup's action parameter to Physlib's standard ℏ
\end{definition}
\begin{proof}
  This is the declared model definition of Uses Standard Reduced Planck Constant.
\end{proof}
\begin{definition}[exterior Schrodinger Left Hand Side]
  \label{def:physics:phyx-mini-0617:phyxminiproblems-problemphyxmini0617-exteriorschrodingerlefthandside}
  \lean{PhyXMiniProblems.ProblemPhyXMini0617.exteriorSchrodingerLeftHandSide}
  \uses{def:physics:phyx-mini-0617:phyxminiproblems-problemphyxmini0617-massinkilograms, def:physics:phyx-mini-0617:phyxminiproblems-problemphyxmini0617-energyinjoules, def:physics:phyx-mini-0617:phyxminiproblems-problemphyxmini0617-actioninjouleseconds, def:physics:phyx-mini-0617:phyxminiproblems-problemphyxmini0617-finitepotentialwellsetup}
  The kinetic-plus-potential side of the stationary one-dimensional Schrödinger equation in the exterior constant-potential region
\end{definition}
\begin{proof}
  This is the declared model definition of exterior Schrodinger Left Hand Side.
\end{proof}
\begin{definition}[Satisfies Exterior Stationary Schrodinger Equation]
  \label{def:physics:phyx-mini-0617:phyxminiproblems-problemphyxmini0617-satisfiesexteriorstationaryschrodingerequation}
  \lean{PhyXMiniProblems.ProblemPhyXMini0617.SatisfiesExteriorStationarySchrodingerEquation}
  \uses{def:physics:phyx-mini-0617:phyxminiproblems-problemphyxmini0617-lengthinmeters, def:physics:phyx-mini-0617:phyxminiproblems-problemphyxmini0617-energyinjoules, def:physics:phyx-mini-0617:phyxminiproblems-problemphyxmini0617-finitepotentialwellsetup, def:physics:phyx-mini-0617:phyxminiproblems-problemphyxmini0617-exteriorschrodingerlefthandside}
  Twice-continuous differentiability and the stationary Schrödinger equation on x > L. This governing-law premise does not assume an exponential form
\end{definition}
\begin{proof}
  This is the declared model definition of Satisfies Exterior Stationary Schrodinger Equation.
\end{proof}
\begin{definition}[Satisfies Finite Boundary Condition At Infinity]
  \label{def:physics:phyx-mini-0617:phyxminiproblems-problemphyxmini0617-satisfiesfiniteboundaryconditionatinfinity}
  \lean{PhyXMiniProblems.ProblemPhyXMini0617.SatisfiesFiniteBoundaryConditionAtInfinity}
  \uses{def:physics:phyx-mini-0617:phyxminiproblems-problemphyxmini0617-lengthinmeters, def:physics:phyx-mini-0617:phyxminiproblems-problemphyxmini0617-finitepotentialwellsetup}
  The stated boundary condition that the wave function remains finite toward positive infinity, represented without assuming a limit or an exponential ansatz: its norm is uniformly bounded beyond some exterior cutoff
\end{definition}
\begin{proof}
  This is the declared model definition of Satisfies Finite Boundary Condition At Infinity.
\end{proof}
% --- Archon declaration topology end ---
% --- Archon physics formalization source end ---
